\subsection{Optimal natural selector weights}
\label{sec:natural-selectors}

This section gives a completed local improvement to the direct
commutative-certificate architecture.  Its global bounded closure has not
been compiled, so the local dimensions below are not a new final
consistency-equation variable count.

\begin{theorem}[Natural interpolation weights]
Let $k\geq1$, and let $G_0,\ldots,G_{k-1}$ be integer polynomials
nonnegative on natural inputs.  Set
\[
 g_k(r)=\prod_{j=0}^{k-1}(r-j),\qquad
 L_i(r)=\prod_{\substack{0\leq j<k\\j\ne i}}(r-j),
\]
and define
\[
 W_i(r)=
 \begin{cases}
 L_i(r),&k-1-i\text{ even},\\
 (r-k)L_i(r),&k-1-i\text{ odd}.
 \end{cases}
\]
Then $g_k$ and every $W_i$ are nonnegative on $\mathbb N$, and
\[
 \exists r\in\mathbb N:\quad
 g_k(r)+\sum_{i=0}^{k-1}W_i(r)G_i=0
 \quad\Longleftrightarrow\quad
 \bigvee_{i=0}^{k-1}G_i=0.
\]
Each $W_i$ has the smallest possible degree among real polynomials
nonnegative on $\mathbb N$, positive at $i$, and zero at every other
integer in $\{0,\ldots,k-1\}$: degree $k-1$ in the even case and $k$
in the odd case.
\end{theorem}
\begin{proof}
For $r<k$, the factors give $g_k(r)=0$ and $W_i(r)=0$ unless $r=i$.
At $r=i$, the sign of $L_i(i)$ is $(-1)^{k-1-i}$; the optional factor
$i-k$ changes the negative sign.  For $r\geq k$, all relevant factors
are nonnegative, and $g_k(r)>0$.  Thus every summand in the displayed
equation is nonnegative, a zero forces $r<k$, and its unique surviving
branch weight is positive.  This proves both directions.

For optimality, $k-1$ required distinct roots force degree at least
$k-1$.  At that degree a candidate equals $cL_i$; nonnegativity at
large natural arguments forces $c>0$.  Its value at $i$ has the required
sign exactly in the even case.  The odd case therefore needs degree at
least $k$, attained by the stated construction.
\end{proof}

The incoming Round11 source uses twelve nonnegative Horn branch
polynomials $h_0,\ldots,h_{11}$, with respective degrees $2$, ten
copies of $4$, and $8$ for the bundle branch $h_9$.  The truth branch is
$h_0=p^2+q^2+v^2$.  Retain $M=M_0+1$ and the accumulator residuals
\begin{align*}
 E&=S'(1+vM)-S(1+pM)(1+qM),\\
 F&=R'-R\bigl((1+p)(1+q)+v\bigr).
\end{align*}
At a truth step $p=q=v=0$, the equations say $S'=S$ and $R'=R$.
All such steps can therefore be deleted from a finite certificate,
preserving its endpoints and final height.  The existential trace length
and its finite CRT code are replaced by those of the shortened trace.
The distinguished endpoints are $S_0=M+1$ and $S_N=1$, with $M>0$, so a successful trace cannot consist entirely of deleted steps. Conversely every trace without these steps is already an old trace.
This establishes equality of the global existential source predicates;
it does not assert equality of their local zero sets at a fixed step.

Group the remaining eleven factors as
\[
 (G_0,\ldots,G_5)=
 (h_9,h_1h_2,h_3h_4,h_5h_6,h_7h_8,h_{10}h_{11}).
\]
Every group has degree eight.  The resulting literal local polynomial is
\[
 \mathcal L=g_6(r)+\sum_{i=0}^{5}W_i(r)G_i+E^2+F^2.
\]
Its degree is exactly $14$, with $21$ local natural witnesses and five
exposed state/parameter coordinates.  The former three-bit construction
uses $23$ local witnesses at the same degree.  Exact sparse expansion of
the supplied rows gives $50\,461$ monomials; the leading term of
$\mathcal L(t,\ldots,t)$ is $723072t^{14}$.  The source JSON, selector
weights, entire expanded polynomial, and a standard-library verifier are
included.  The exact degree assertion is therefore independently
reproducible.

There is a second improvement using two ternary selectors $u,w$.  The
$k=3$ weights are
\[
 U_0(x)=(x-1)(x-2),\quad
 U_1(x)=x(x-2)(x-3),\quad
 U_2(x)=x(x-1),
\]
of degrees $2,3,2$.  Use the group array
\[
 (G_{ij})=
 \begin{pmatrix}
 h_9&h_5&h_1h_2\\
 h_6&h_7&h_8\\
 h_3h_4&h_{10}&h_{11}
 \end{pmatrix}
\]
and the polynomial
\[
 \mathcal L_{12}=g_3(u)+g_3(w)
 +\sum_{i,j=0}^{2}U_i(u)U_j(w)G_{ij}+E^2+F^2.
\]
The degree-eight groups occupy corners with weight degree four.  Every
other group contributes at degree at most ten.  Thus this source has
$22$ local witnesses and degree exactly $12$; its $40\,990$-monomial
expansion has leading coefficient $627840$ under uniform substitution.
Nonnegativity and the two selector penalties prove the same local and
global source equivalence.  The new local tradeoffs are consequently
$(21,14)$ and $(22,12)$ in place of $(23,14)$.

The incoming index-free inert-prime closure has conditional coordinate
count $2m+9+2f+p+(m+1)b$, for local witness dimension $m$ and chosen
factorial, power, and binomial graph costs $f,p,b$.  Holding the remaining
graph conventions fixed, the change $m=23\longrightarrow21$ saves
$4+2b$ coordinates at that interface.  This is a conditional compiler
accounting identity; it does not replace materialization of the entire
direct consistency polynomial.  The degree-twelve alternative instead
saves $2+b$ coordinates under that same convention.

\subsection{Exact local arithmetic graphs}

A second improvement exploits the natural domain directly.  For natural
$x_1,\ldots,x_s$, their sum vanishes exactly when they all vanish.
Hence a mask $(\sum x_i)F^2$ can replace $(\sum x_i^2)F^2$ without
changing its zeros.  In particular,
\[
 C(n,p)=(n+p)(n-p-1)^2
\]
is a nonnegative cubic with zeros exactly at $p=\operatorname{pred}(n)$.
The corresponding signed quadratic residual $(n+p)(n-p-1)$ has the same
zero set and may be squared when a strict residual-system compiler is
required.

More significantly, the quadratic polynomial
\[
 Z(n,b,p)=(n+b-p-1)^2+bp
\]
jointly represents predecessor and zero test:
\[
 Z=0\quad\Longleftrightarrow\quad
 p=\operatorname{pred}(n),\qquad b=\mathbf1_{n=0}.
\]
Indeed $bp=0$ and $n+b=p+1$.  If $b=0$, then $n=p+1\geq1$; if
$p=0$, then $n+b=1$, giving only $(n,b)=(0,1)$ or $(1,0)$.
The converse is immediate.  Hiding $p$ gives a one-witness zero-test
graph of degree two.  If both outputs are already present, no extra
witness is needed.

For a conditional, one may use the one-witness quartic
\[
 n(c-b)^2+\bigl((c-a)(n-s-1)\bigr)^2,
\]
or the two-witness cubic
\[
 Z(n,t,p)+t(c-a)^2+n(c-b)^2.
\]
Their zeros project exactly to $c=a$ when $n=0$, and $c=b$ otherwise.
For the quartic, $n=0$ makes $n-s-1$ nonzero for every natural $s$;
$n>0$ forces the first branch directly.  For the cubic, the joint graph
fixes the exact branch bit.  Completeness uses the canonical predecessor
and bit.

Keeping the incoming two-witness max-square pairing polynomial $Q$
unchanged, the tree projections use $Q(p,a,b;t,s)+C(n,p)$.  The actual
found-marker list head uses
\[
 Q(p,a,b;t,s)+(n+p+c)(n-p-1)^2+n(c-a-1)^2.
\]
At $n=0$ the middle term forces $p=c=0$; at $n>0$ it forces $p=n-1$
and the last term forces $c=a+1$.  Thus nil and the stored-zero marker
are both preserved.  All naturalness assumptions are explicit.

\begin{center}
\begin{tabular}{lccc}
\hline
Local graph & Hidden witnesses & Degree & Tree operations\\
\hline
Predecessor & $0$ & $4\to3$ & $10\to8$\\
Zero test & $1$ & $4\to2$ & $14\to10$\\
Tree projection/list tail & $4$ & $4$ & $49\to47$\\
Conditional & $1$ & $4$ & $16\to15$\\
List head & $5$ & $4$ & $59\to55$\\
\hline
\end{tabular}
\end{center}

The comparison is against the incoming Joint Guards package and counts
literal binary-expression trees, including repeated square operands and
the same final addition of zero.  Full expansions and separate shared-DAG
counts are supplied.  These are exact local graph improvements, without
a claimed end-to-end compiler or object-level formal certificate.
